\chapter{Complexity — Formula Fan Out}
%%% generated by the blueprint pipeline from TCSlib.Complexity.CircuitComplexity.FormulaFanOut
\section{Overview}

Declarations of the Lean module \texttt{TCSlib.Complexity.CircuitComplexity.FormulaFanOut}.

\section{Declarations}

\begin{definition}[Single-root fan-out conditions]
\lean{BoolCircuit.CompileFanout}
\label{BoolCircuit.CompileFanout}
\leanok
\uses{BoolCircuit.DAGGate, BoolCircuit.fanoutIn}
Let $n$ be a natural number, let $G$ and $E$ be finite gate lists, and let $r$ be a
natural number. The single-root fan-out conditions require that every old vertex $v$
with $n\le v<n+|G|$ has fan-out zero in $E$, and every vertex $v\ge n$ has fan-out at
most one in $E$. If $r\ge n$, then $r$ has fan-out zero; moreover, either $r<n$ or $r\ge
n+|G|$. Every vertex $v$ with $n+|G|\le v<n+|G|+|E|$ other than $r$ has fan-out exactly
one in $E$.
\end{definition}

\begin{definition}[Multiple-root fan-out conditions]
\lean{BoolCircuit.CompileListFanout}
\label{BoolCircuit.CompileListFanout}
\leanok
\uses{BoolCircuit.DAGGate, BoolCircuit.fanoutIn}
Let $n$ be a natural number, let $G$ and $E$ be finite gate lists, and let $R$ be a list
of natural-number roots. Every old vertex $v$ with $n\le v<n+|G|$ has fan-out zero in
$E$, and every vertex $v\ge n$ has fan-out at most one in $E$. Each root $r$ in $R$ with
$r\ge n$ has fan-out zero, and every root is either below $n$ or at least $n+|G|$. Every
vertex $v$ with $n+|G|\le v<n+|G|+|E|$ that is not in $R$ has fan-out exactly one in
$E$.
\end{definition}

\begin{theorem}[Zero fanout beyond the circuit]
\lean{BoolCircuit.fanoutIn_ext_eq_zero}
\label{BoolCircuit.fanoutIn_ext_eq_zero}
\leanok
\uses{BoolCircuit.DAGGate, BoolCircuit.GatesAcyclic, BoolCircuit.fanoutIn, BoolCircuit.fanoutIn_append, BoolCircuit.fanoutIn_eq_zero_of_le}
\difficulty{3}
Let $n$ be a nonnegative integer, and let $gs$ and $ext$ be lists of Boolean gates such
that every gate in their concatenation reads only vertices with index less than $n+i$,
where $i$ is that gate’s position in the concatenation. If $v$ satisfies $n+|gs++ext|\le
v$, then no gate in $ext$ reads vertex $v$; equivalently, the fanout of $v$ within $ext$
is zero.
\end{theorem}

\begin{theorem}[Fanout of compiled circuit lists]
\lean{BoolCircuit.compileTrees_fanout}
\label{BoolCircuit.compileTrees_fanout}
\leanok
\uses{BoolCircuit.CompileFanout, BoolCircuit.CompileListFanout, BoolCircuit.GatesAcyclic, BoolCircuit.TreeCircuit, BoolCircuit.compileTree, BoolCircuit.compileTree_spec, BoolCircuit.compileTrees, BoolCircuit.compileTrees_spec, BoolCircuit.fanoutIn, BoolCircuit.fanoutIn_append, BoolCircuit.fanoutIn_ext_eq_zero}
\difficulty{6}
Let $n$ be a natural number and $cs$ a finite list of Boolean circuit trees on $n$
variables. Suppose that, for every $c$ in $cs$, compiling $c$ after any acyclic gate
list $gs$ appends a list of gates satisfying the single-circuit fanout conditions: old
gate vertices are unused, every vertex of index at least $n$ is used at most once, any
root of index at least $n$ is unused, each root is either an input vertex of index less
than $n$ or a new vertex beyond the old gates, and every new vertex other than the root
is used exactly once. Then for every acyclic initial gate list $gs$, compiling the whole
list $cs$ appends some gate list $ext$ such that old gate vertices are unused, each
vertex of index at least $n$ is used at most once, every listed root of index at least
$n$ is unused, each listed root is either an input vertex of index less than $n$ or a
new vertex beyond the old gates, and every new vertex not among the listed roots is used
exactly once.
\end{theorem}

\begin{theorem}[Fanout of a compiled circuit]
\lean{BoolCircuit.compileTree_fanout}
\label{BoolCircuit.compileTree_fanout}
\leanok
\uses{BoolCircuit.CompileFanout, BoolCircuit.DAGGate, BoolCircuit.GatesAcyclic, BoolCircuit.TreeCircuit, BoolCircuit.TreeCircuit.ind, BoolCircuit.compileTree, BoolCircuit.compileTree_spec, BoolCircuit.compileTrees, BoolCircuit.compileTrees_fanout, BoolCircuit.compileTrees_spec, BoolCircuit.fanoutIn, BoolCircuit.fanoutIn_append, BoolCircuit.fanoutIn_ext_eq_zero, BoolCircuit.fanoutIn_singleton, BoolCircuit.nodeKind}
\difficulty{6}
Let $n$ be a natural number, let $c$ be a Boolean circuit tree on $n$ variables, and let
$gs$ be a finite list of gates on natural-number-indexed vertices such that every input
to its $i$th gate has index less than $n+i$. Compiling $c$ after $gs$ appends some list
$ext$ of gates. Among the appended gates, each vertex of index at least $n$ is used as
an input at most once; vertices corresponding to gates in $gs$ are not used, and the
root vertex is not used when its index is at least $n$. The root index is either less
than $n$ or at least $n+|gs|$. Every vertex whose index is at least $n+|gs|$ and less
than $n+|gs|+|ext|$, except the root, is used exactly once as an input to the appended
gates.
\end{theorem}

\begin{theorem}[Fanout conditions for the circuit DAG translation]
\lean{BoolCircuit.TreeCircuit.toDAG_fanoutSpec}
\label{BoolCircuit.TreeCircuit.toDAG_fanoutSpec}
\leanok
\uses{BoolCircuit.CompileFanout, BoolCircuit.GatesAcyclic.nil, BoolCircuit.TreeCircuit, BoolCircuit.TreeCircuit.toDAG, BoolCircuit.compileTree_fanout}
\difficulty{3}
For every Boolean circuit $c$ on $n$ variables, let $G$ be the list of gates in its
directed acyclic graph translation, let $r$ be its output index, and let $F(v)$ count
the gates in $G$ whose argument list contains the index $v$. Every index $v$ with $n\le
v<n+|G|$ has $F(v)=0$, every $v\ge n$ has $F(v)\le1$, and if $r\ge n$ then $F(r)=0$.
Also, either $r<n$ or $n+|G|\le r$, and every index $v$ with $n+|G|\le v<n+|G|+|G|$ and
$v\ne r$ has $F(v)=1$.
\end{theorem}

\begin{theorem}[Fanout bound for compiled circuits]
\lean{BoolCircuit.TreeCircuit.toDAG_fanout_le_one}
\label{BoolCircuit.TreeCircuit.toDAG_fanout_le_one}
\leanok
\uses{BoolCircuit.DAGCircuit.fanout, BoolCircuit.TreeCircuit, BoolCircuit.TreeCircuit.toDAG, BoolCircuit.TreeCircuit.toDAG_fanoutSpec}
\difficulty{1}
Let $c$ be a Boolean circuit tree on $n$ variables, and compile it into a directed
acyclic circuit. Every vertex $v \ge n$ is read by at most one gate.
\end{theorem}

\begin{theorem}[Output fanout of compiled circuits]
\lean{BoolCircuit.TreeCircuit.toDAG_fanout_output}
\label{BoolCircuit.TreeCircuit.toDAG_fanout_output}
\leanok
\uses{BoolCircuit.DAGCircuit.fanout, BoolCircuit.TreeCircuit, BoolCircuit.TreeCircuit.toDAG, BoolCircuit.TreeCircuit.toDAG_fanoutSpec}
\difficulty{1}
Let $c$ be a Boolean circuit tree on $n$ variables, and compile it into a directed
acyclic circuit. If the output vertex is at least $n$, then no gate reads the output
vertex.
\end{theorem}

\begin{theorem}[Fanout of non-output gates]
\lean{BoolCircuit.TreeCircuit.toDAG_fanout_eq_one}
\label{BoolCircuit.TreeCircuit.toDAG_fanout_eq_one}
\leanok
\uses{BoolCircuit.DAGCircuit.fanout, BoolCircuit.DAGCircuit.size, BoolCircuit.TreeCircuit, BoolCircuit.TreeCircuit.toDAG, BoolCircuit.TreeCircuit.toDAG_fanoutSpec}
\difficulty{1}
Let $c$ be a Boolean circuit tree on $n$ variables, and compile it into a directed
acyclic circuit. For any vertex $v$ with $n \le v < \mathrm{size}(c\text{ compiled})$
that is not the output vertex, the number of gates reading $v$ is exactly $1$.
\end{theorem}

\begin{theorem}[Repeated input use in a converted circuit]
\lean{BoolCircuit.TreeCircuit.exists_toDAG_input_fanout_two}
\label{BoolCircuit.TreeCircuit.exists_toDAG_input_fanout_two}
\leanok
\uses{BoolCircuit.DAGCircuit.fanout, BoolCircuit.TreeCircuit, BoolCircuit.TreeCircuit.toDAG, BoolCircuit.compileTree, BoolCircuit.compileTrees, BoolCircuit.fanoutIn, BoolCircuit.nodeKind}
\difficulty{1}
There exists a Boolean circuit tree on one variable whose associated directed acyclic
circuit has fanout exactly $2$ at input vertex $0$.
\end{theorem}

\begin{theorem}[Bounded-fanin tree compilation]
\lean{BoolCircuit.TreeCircuit.exists_dagCircuit_fanout_le_one}
\label{BoolCircuit.TreeCircuit.exists_dagCircuit_fanout_le_one}
\leanok
\uses{BoolCircuit.DAGCircuit, BoolCircuit.DAGCircuit.IsFaninTwo, BoolCircuit.DAGCircuit.eval, BoolCircuit.DAGCircuit.fanout, BoolCircuit.DAGCircuit.size, BoolCircuit.TreeCircuit, BoolCircuit.TreeCircuit.eval, BoolCircuit.TreeCircuit.maxFanin, BoolCircuit.TreeCircuit.size, BoolCircuit.TreeCircuit.toDAG, BoolCircuit.TreeCircuit.toDAG_eval, BoolCircuit.TreeCircuit.toDAG_fanout_le_one, BoolCircuit.TreeCircuit.toDAG_isFaninTwo, BoolCircuit.TreeCircuit.toDAG_size_le}
\difficulty{1}
Let $c$ be a Boolean circuit tree on $n$ variables whose maximum fanin is at most $2$.
There exists an equivalent acyclic Boolean circuit with gates of fanin at most $2$, with
no repeated input to a gate and exactly one input to each NOT gate, such that every
non-input vertex is read by at most one gate and the circuit size is at most
$n+\mathrm{size}(c)$.
\end{theorem}

\begin{theorem}[Congruence of gate unfolding]
\lean{BoolCircuit.unfoldGate_congr}
\label{BoolCircuit.unfoldGate_congr}
\leanok
\uses{BoolCircuit.DAGGate, BoolCircuit.DAGGate.WellFormed, BoolCircuit.TreeCircuit, BoolCircuit.unfoldGate}
\difficulty{4}
Let $n$ be a natural number, let $g$ be a well-formed gate, and let $t,t'$ assign a
Boolean circuit on $n$ variables to each gate argument and polarity. If $t(a,q)=t'(a,q)$
for every argument $a$ of $g$ and every polarity $q$, then unfolding $g$ using either
assignment gives the same circuit at every polarity $p$.
\end{theorem}

\begin{theorem}[Unfolded gate size is polarity-independent]
\lean{BoolCircuit.size_unfoldGate_polarity}
\label{BoolCircuit.size_unfoldGate_polarity}
\leanok
\uses{BoolCircuit.DAGGate, BoolCircuit.TreeCircuit, BoolCircuit.TreeCircuit.size, BoolCircuit.TreeCircuit.size_node, BoolCircuit.sumSize_map, BoolCircuit.unfoldGate}
\difficulty{4}
Let $g$ be any gate, and let $t$ assign a Boolean circuit tree to each vertex and
polarity. Suppose that at every vertex the assigned tree has the same size for both
polarities. Then unfolding $g$ at either of two polarities gives trees of equal size.
\end{theorem}

\begin{theorem}[Size bound for an unfolded gate]
\lean{BoolCircuit.size_unfoldGate_le}
\label{BoolCircuit.size_unfoldGate_le}
\leanok
\uses{BoolCircuit.DAGGate, BoolCircuit.DAGGate.WellFormed, BoolCircuit.TreeCircuit, BoolCircuit.TreeCircuit.size, BoolCircuit.TreeCircuit.size_node, BoolCircuit.sumSize_map, BoolCircuit.unfoldGate}
\difficulty{4}
Let $g$ be a gate whose argument list has no repetitions and, if $g$ is a NOT gate, has
length one. Let $t$ assign a Boolean circuit tree to each vertex and polarity, and
suppose that for every vertex the assigned circuit has the same size at both polarities.
For either polarity $p$, the size of the circuit obtained by unfolding $g$ at $p$ is at
most one plus the sum of the sizes of the circuits assigned to its arguments at polarity
$p$.
\end{theorem}

\begin{theorem}[Unfolded vertex size is polarity-independent]
\lean{BoolCircuit.size_unfoldVertex_polarity}
\label{BoolCircuit.size_unfoldVertex_polarity}
\leanok
\uses{BoolCircuit.DAGGate, BoolCircuit.TreeCircuit.size, BoolCircuit.size_unfoldGate_polarity, BoolCircuit.unfoldVertex}
\difficulty{4}
Fix a natural number $n$ and a finite list of gates. For every fuel bound, vertex, and
pair of Boolean polarities, the circuit obtained by unfolding that vertex with the first
polarity has the same size as the circuit obtained with the second polarity.
\end{theorem}

\begin{theorem}[Fuel independence of circuit unfolding]
\lean{BoolCircuit.DAGCircuit.unfoldVertex_fuel_congr}
\label{BoolCircuit.DAGCircuit.unfoldVertex_fuel_congr}
\leanok
\uses{BoolCircuit.DAGCircuit.IsWellFormed, BoolCircuit.constGate, BoolCircuit.unfoldGate, BoolCircuit.unfoldGate_congr, BoolCircuit.unfoldVertex}
\difficulty{5}
Let $C$ be a Boolean directed acyclic circuit with $n$ inputs, whose gates have distinct
argument vertices and whose NOT gates each have exactly one argument. For any vertex $v$
and nonnegative fuel bounds $f,f'$ with $v<f$ and $v<f'$, unfolding $v$ with either fuel
bound and the same polarity produces the same Boolean circuit tree.
\end{theorem}

\begin{definition}[Size of a vertex unfolding]
\lean{BoolCircuit.DAGCircuit.unfoldSize}
\label{BoolCircuit.DAGCircuit.unfoldSize}
\leanok
\uses{BoolCircuit.TreeCircuit.size, BoolCircuit.unfoldVertex}
Let $n$ be a natural number and let $C$ be a directed acyclic circuit with $n$ inputs.
For each natural-number vertex $v$, define its unfolding size to be the number of
entries in the list produced by unfolding $v$ in $C$, with bound $v+1$ and the initial
flag set to true.
\end{definition}

\begin{theorem}[Unfolded size at an input]
\lean{BoolCircuit.DAGCircuit.unfoldSize_input}
\label{BoolCircuit.DAGCircuit.unfoldSize_input}
\leanok
\uses{BoolCircuit.DAGCircuit.unfoldSize, BoolCircuit.TreeCircuit.size, BoolCircuit.unfoldVertex}
\difficulty{1}
Let $C$ be a Boolean directed acyclic circuit with $n$ input variables. For every vertex
index $v<n$, the size obtained by unfolding that input vertex is $1$.
\end{theorem}

\begin{theorem}[Unfolded size at a gate]
\lean{BoolCircuit.DAGCircuit.unfoldSize_gate_le}
\label{BoolCircuit.DAGCircuit.unfoldSize_gate_le}
\leanok
\uses{BoolCircuit.DAGCircuit.IsFaninTwo, BoolCircuit.DAGCircuit.unfoldSize, BoolCircuit.DAGCircuit.unfoldSize_input, BoolCircuit.DAGCircuit.unfoldVertex_fuel_congr, BoolCircuit.DAGGate.WellFormed, BoolCircuit.TreeCircuit.size, BoolCircuit.size_unfoldGate_le, BoolCircuit.size_unfoldVertex_polarity, BoolCircuit.unfoldGate, BoolCircuit.unfoldVertex}
\difficulty{5}
Let $C$ be a Boolean directed acyclic circuit whose gates have pairwise distinct input
vertices, whose NOT gates each have exactly one input, and whose gates each have at most
two input vertices. For any gate index $k$ within the gate list, the size obtained by
unfolding the vertex corresponding to gate $k$ is at most $3$ plus the sum of the
unfolded sizes of the distinct gate-input vertices of that gate (those with index at
least the number of input variables).
\end{theorem}

\begin{definition}[Roots among an initial segment of gates]
\lean{BoolCircuit.DAGCircuit.rootSet}
\label{BoolCircuit.DAGCircuit.rootSet}
\leanok
\uses{BoolCircuit.fanoutIn}
Let $C$ be a finite acyclic Boolean circuit with $n$ input vertices, and let $k$ be a
natural number. The root set consists of the gate vertices with indices in the half-open
interval $[n,n+k)$ that have fan-out zero among the first $k$ gates of $C$.
\end{definition}

\begin{theorem}[Acyclicity of a gate prefix]
\lean{BoolCircuit.DAGCircuit.gatesAcyclic_take}
\label{BoolCircuit.DAGCircuit.gatesAcyclic_take}
\leanok
\uses{BoolCircuit.GatesAcyclic, BoolCircuit.gatesAcyclic_append}
\difficulty{1}
Let $C$ be a finite directed acyclic Boolean circuit with $n$ input vertices, and let
$k$ be a natural number. The list consisting of the first $k$ gates of $C$ is acyclic:
each argument of its gate at position $i$ is a vertex with index less than $n+i$.
\end{theorem}

\begin{theorem}[Successive root sets]
\lean{BoolCircuit.DAGCircuit.rootSet_succ}
\label{BoolCircuit.DAGCircuit.rootSet_succ}
\leanok
\uses{BoolCircuit.DAGCircuit.gatesAcyclic_take, BoolCircuit.DAGCircuit.rootSet, BoolCircuit.fanoutIn, BoolCircuit.fanoutIn_append, BoolCircuit.fanoutIn_eq_zero_of_le, BoolCircuit.fanoutIn_singleton}
\difficulty{5}
Let $C$ be a directed acyclic circuit with $n$ input vertices, and let $k$ be a natural
number smaller than the number of gates in $C$. The set of vertices in the interval
$[n,n+k+1)$ that are read by none of the first $k+1$ gates is obtained by adjoining the
vertex $n+k$ to the vertices in the corresponding set for the first $k$ gates that are
not arguments of gate $k$.
\end{theorem}

\begin{theorem}[Bound on unfolding root sizes]
\lean{BoolCircuit.DAGCircuit.sum_rootSet_le}
\label{BoolCircuit.DAGCircuit.sum_rootSet_le}
\leanok
\uses{BoolCircuit.DAGCircuit.IsFaninTwo, BoolCircuit.DAGCircuit.fanout, BoolCircuit.DAGCircuit.rootSet, BoolCircuit.DAGCircuit.rootSet_succ, BoolCircuit.DAGCircuit.unfoldSize, BoolCircuit.DAGCircuit.unfoldSize_gate_le, BoolCircuit.fanoutIn, BoolCircuit.fanoutIn_append, BoolCircuit.fanoutIn_singleton}
\difficulty{6}
Let $C$ be a directed acyclic Boolean circuit on $n$ inputs in which each gate has
distinct arguments, every NOT gate has exactly one argument, and every gate has at most
two arguments. Suppose also that every vertex with index at least $n$ is an argument of
at most one gate. For every $k$ with $k\leq$ the number of gates, the sum of the
unfolded sizes over the vertices in $[n,n+k)$ that have fanout zero among the first $k$
gates is at most $3k$.
\end{theorem}

\begin{theorem}[Unfolded size with bounded fanout]
\lean{BoolCircuit.DAGCircuit.toTree_size_le_of_fanout_le_one}
\label{BoolCircuit.DAGCircuit.toTree_size_le_of_fanout_le_one}
\leanok
\uses{BoolCircuit.DAGCircuit.IsFaninTwo, BoolCircuit.DAGCircuit.fanout, BoolCircuit.DAGCircuit.rootSet, BoolCircuit.DAGCircuit.rootSet_succ, BoolCircuit.DAGCircuit.sum_rootSet_le, BoolCircuit.DAGCircuit.toTree, BoolCircuit.DAGCircuit.unfoldSize, BoolCircuit.DAGCircuit.unfoldSize_input, BoolCircuit.TreeCircuit.size}
\difficulty{4}
Let $C$ be a Boolean directed acyclic circuit whose gates have pairwise distinct input
vertices, whose NOT gates each have exactly one input, and whose gates each have at most
two input vertices. Suppose every non-input vertex is used as an input by at most one
gate. Then the size of the tree obtained by unfolding $C$ is at most $3$ times the
number of gates of $C$, plus $1$.
\end{theorem}

\begin{theorem}[Tree representation under low fanout]
\lean{BoolCircuit.DAGCircuit.exists_treeCircuit_of_fanout_le_one}
\label{BoolCircuit.DAGCircuit.exists_treeCircuit_of_fanout_le_one}
\leanok
\uses{BoolCircuit.DAGCircuit.IsFaninTwo, BoolCircuit.DAGCircuit.eval, BoolCircuit.DAGCircuit.fanout, BoolCircuit.DAGCircuit.size, BoolCircuit.DAGCircuit.toTree, BoolCircuit.DAGCircuit.toTree_eval, BoolCircuit.DAGCircuit.toTree_maxFanin_le, BoolCircuit.DAGCircuit.toTree_size_le_of_fanout_le_one, BoolCircuit.TreeCircuit, BoolCircuit.TreeCircuit.eval, BoolCircuit.TreeCircuit.maxFanin, BoolCircuit.TreeCircuit.size}
\difficulty{3}
Let $n$ be a nonnegative integer, and let $C$ be a directed acyclic Boolean circuit on
$n$ inputs in which every gate has distinct arguments and at most two arguments, and
every NOT gate has exactly one argument. Suppose every vertex with index at least $n$
has fanout at most $1$. Then there is a Boolean circuit tree $t$ on the same inputs that
computes the same function as $C$, has maximum fanin at most $2$, and satisfies both
size bounds $|t|\leq 3\,\mathrm{length}(C.\mathrm{gates})+1$ and $|t|\leq
4\,\mathrm{size}(C)$.
\end{theorem}
